\documentclass[11pt]{article}
\usepackage{setspace}
\usepackage{geometry}
\geometry{margin=1in}
\setstretch{1.15}

\begin{document}

\noindent
\textbf{Zhongli Wang} \\
Phone: +49 160 9541 6015 \\
Email: nkwangzhongli@gmail.com \\
Website: https://sites.google.com/view/zhongli-wang \\
Bielefeld, Germany \\

\vspace{0.45cm}

\noindent
\textbf{Selection Committee} \\
School of Law, College of Social Sciences \\
University of Glasgow \\

\vspace{0.45cm}

\noindent
\textbf{Application for Research Associate --- Civic Finance (ERC Starting Grant \#101165302)}\\

\vspace{0.25cm}

\noindent Dear Professor Solana and Members of the Selection Committee,\\[0.15cm]

I am writing to apply for the Research Associate position on \emph{Civic Finance} at the University of Glasgow. I am currently completing my Ph.D.\ in Economics at Bielefeld University and Universit\`a Cattolica del Sacro Cuore (Marie Sk{\l}odowska--Curie EPOC programme). My research focuses on complexity in economic systems---how heterogeneous agents interact through networks, how incentives and information shape strategic behavior, and how institutional interventions change system-level outcomes. The Civic Finance agenda links strongly with my own interests: building empirically grounded, complexity-informed models that can explain and improve real-world outcomes.
% ---here, the financial services that enable high-emitting activity and the legal levers that could steer finance toward a civic transition.

My doctoral work develops agent-based models of market environments where outcomes are non-linear, path-dependent, and shaped by feedback loops. In my dissertation, \emph{Strategic Behavior and Economic Dynamics in ICT Markets: An Agent-Based Framework}, I study how firms and heterogeneous users interact under network effects, congestion constraints, and imperfect information. Methodologically, this work required (i) formalizing behavioral and strategic decision rules, (ii) building scalable simulation pipelines, and (iii) conducting counterfactual analysis to evaluate policy-relevant interventions. These skills translate directly to Civic Finance: the project’s aim---to understand how financial institutions and high-emitting actors interact under legal-institutional constraints---is naturally suited to an agent-based approach that can represent strategic adaptation, regulatory avoidance, and unintended spillovers.

A central contribution I can make to Civic Finance is to connect three components into a coherent research pipeline: (1) data construction on financial services and their links to high-emitting actors, (2) network identification of influential institutions and key channels of financial dependence, and (3) an agent-based model that can stress-test legal interventions under realistic behavioral responses. I have substantial experience operationalizing network concepts in economic settings: defining nodes and edges from messy sources, constructing interpretable measures of influence and centrality, and translating network structure into behavioral mechanisms in simulations. In Civic Finance, this would support the identification of financial institutions whose underwriting, lending, asset management, or advisory services are most central for sustaining high-emitting activity, and would provide a rigorous basis for prioritizing regulatory focus.

I am comfortable working with large, heterogeneous datasets and building reproducible data infrastructure. In collaborative work on \emph{Polytope Fraud Theory}, I used supervised machine learning to detect structured anomalies in financial statements (with a focus on firms targeted by activist short sellers), which strengthened my experience in feature engineering, robustness checks, and translating algorithmic signals into economic interpretation. In a second project, \emph{Agentic Workflows for Economic Research}, I helped develop a framework for integrating large-language-model agents into empirical pipelines (e.g., document collection and classification, transparent workflow documentation, etc.). For Civic Finance, these experiences are directly relevant for constructing and maintaining the project database, documenting extraction and coding procedures, and ensuring that data and results remain auditable and reusable for papers and policy outputs.

If selected, my near-term plan would be to deliver publishable outputs through a staged strategy aligned with the project’s deliverables. First, I would help build the core dataset linking financial institutions to high-emitting market actors and categorizing the relevant services, with careful attention to measurement validity, legal definitions, and reproducibility (including codebooks, versioning, and transparent data provenance). Second, I would implement network analyses to identify influential institutions and channels (e.g., syndication structures, repeated relationships, and concentration of service provision), producing a first co-authored manuscript focused on mapping and measurement. Third, I would develop an agent-based model that combines institution heterogeneity (business models, risk constraints, and compliance incentives) with actor heterogeneity (financing needs, emission profiles, and substitution options) to evaluate the effects of alternative legal interventions. The model would be designed to produce testable implications and clear policy counterfactuals, supporting a second manuscript centered on mechanisms and intervention design.

Beyond technical fit, I am motivated by the project’s interdisciplinary setting and the opportunity to work at the boundary of law, finance, and complex systems. I would actively participate in the Corporate \& Financial Law and Just Transitions clusters and contribute to seminars and workshops. I also value collaborative team practices and can take initiative in coordination tasks, mentoring research assistants or junior team members, and supporting project reporting and budget-related administration when needed.

% The University of Glasgow’s commitment to researcher development is also attractive. I am seeking a role where I can deliver high-quality publications while building a professional profile at the intersection of applied modelling, data infrastructure, and policy-relevant research. The prospect of developing a bespoke development programme with guidance from Professor Solana and the Research Support Office, and of building a global network through the project’s advisory and practitioner links, fits well with my medium-term goal of advancing an internationally visible research agenda on the governance of complex socio-economic systems.

Thank you so much for considering my application! I would welcome the opportunity to discuss how my experience in network analysis, agent-based modeling, and reproducible data-driven research can contribute to Civic Finance’s mission and to the University of Glasgow’s wider research community.

\vspace{0.5cm}

\noindent Sincerely,\\[0.9cm]
\noindent Zhongli Wang

\end{document}